\section{Prime Numbers and Unique Factorisation}

\begin{definition}[Prime number]\label{def:prime}
    An integer $p \ge 2$ is prime if its only positive divisors are $1$ and $p$. An integer $n \ge 2$ that is not prime is called composite.
\end{definition}

\begin{lemma}\label{lem:prime_divisor}
    Every integer $n \ge 2$ has a prime divisor. If $n$ is composite, it has a prime divisor $p \le \sqrt{n}$.
\end{lemma}
\begin{proof}
    Let $p$ be the smallest divisor of $n$ with $p \ge 2$. A divisor $d$ of $p$ with $1 < d < p$ would divide $n$ and contradict minimality, so $p$ is prime. If $n$ is composite, write $n = pm$ with $m \ge 2$. Then $m$ is a divisor of $n$ with $m \ge 2$, so $m \ge p$ and $n = pm \ge p^2$.
\end{proof}
The second statement is the principle of the sieve of Eratosthenes: to test whether $n$ is prime, it is enough to try the primes up to $\sqrt{n}$.

\begin{theorem}[Gauss's lemma]\label{thm:gauss_lemma}
    Let $a,b,c \in \mathbb{Z}$. If $a \mid bc$ and $(a,b) = 1$, then $a \mid c$.
\end{theorem}
\begin{proof}
    By \cref{cor:coprime_bezout}, $ax + by = 1$ for some $x,y$. Then $c = acx + bcy$, and $a$ divides both terms.
\end{proof}

\begin{corollary}[Euclid's lemma]\label{cor:euclid_lemma}
    If $p$ is prime and $p \mid ab$, then $p \mid a$ or $p \mid b$. More generally, if $p \mid a_1 \cdots a_k$, then $p \mid a_i$ for some $i$.
\end{corollary}
\begin{proof}
    The only positive divisors of $p$ are $1$ and $p$, so $(p,a) \in \{1, p\}$. If $p \nmid a$, then $(p,a) = 1$ and \cref{thm:gauss_lemma} gives $p \mid b$. The general case follows by induction on $k$.
\end{proof}

\begin{theorem}[Fundamental theorem of arithmetic]\label{thm:fta}
    Every integer $n \ge 2$ is a product of primes, $n = p_1 p_2 \cdots p_k$, and this factorisation is unique up to the order of the factors.
\end{theorem}
\begin{proof}
    \emph{Existence.} We use strong induction on $n$. If $n$ is prime there is nothing to prove. Otherwise $n = ab$ with $1 < a, b < n$, and the induction hypothesis factors $a$ and $b$.

    \emph{Uniqueness.} Suppose $p_1 \cdots p_k = q_1 \cdots q_l$ with all $p_i, q_j$ prime. Then $p_1$ divides $q_1 \cdots q_l$, so by \cref{cor:euclid_lemma} $p_1 \mid q_j$ for some $j$, hence $p_1 = q_j$. Cancelling this factor and inducting on $k$ gives $k = l$ and the $q_j$ are a rearrangement of the $p_i$.
\end{proof}

Grouping equal primes, every $n \ge 1$ can be written uniquely as
\begin{equation}\label{eq:canonical_factorisation}
    n = \prod_{p} p^{v_p(n)},
\end{equation}
where only finitely many exponents $v_p(n) \ge 0$ are non-zero. The exponent $v_p(n)$ is the $p$-adic valuation of $n$. Uniqueness of the factorisation gives $v_p(mn) = v_p(m) + v_p(n)$.

\begin{corollary}\label{cor:divisors_valuations}
    Let $a,b \ge 1$.
    \begin{enumerate}[label={\upshape\roman*)}]
        \item $a \mid b$ if and only if $v_p(a) \le v_p(b)$ for every prime $p$.
        \item $(a,b) = \prod_p p^{\min(v_p(a), v_p(b))}$ and $[a,b] = \prod_p p^{\max(v_p(a), v_p(b))}$.
        \item $(a,b)[a,b] = ab$.
    \end{enumerate}
\end{corollary}
\begin{proof}
    \begin{enumerate}[label={\upshape\roman*)}]
        \item If $b = ac$, then $v_p(b) = v_p(a) + v_p(c) \ge v_p(a)$. Conversely, if $v_p(a) \le v_p(b)$ for all $p$, then $c = \prod_p p^{v_p(b) - v_p(a)}$ is an integer with $ac = b$.
        \item By i), the common divisors of $a$ and $b$ are the $d$ with $v_p(d) \le \min(v_p(a), v_p(b))$ for all $p$, and the largest of them has equality everywhere. The argument for $[a,b]$ is the same.
        \item Use $\min(x,y) + \max(x,y) = x + y$ in each exponent.
    \end{enumerate}
\end{proof}

\begin{theorem}[Euclid]\label{thm:euclid_infinitude}
    There are infinitely many prime numbers.
\end{theorem}
\begin{proof}
    Let $p_1, \dots, p_k$ be any finite list of primes and set $N = p_1 \cdots p_k + 1$. By \cref{lem:prime_divisor}, $N$ has a prime divisor $p$. If $p = p_i$ for some $i$, then $p$ divides $N - p_1 \cdots p_k = 1$, which is impossible. So $p$ is a prime that is not in the list.
\end{proof}
